% =========================================================
% Compléments M1 — Heston (C. Dérivations)
% =========================================================

\subsection*{(1) Réécrire la dynamique en log-prix : étape indispensable}
Sous $\Qmeas$, Heston est :
\[
\frac{\dd S_t}{S_t}=(r-q)\dd t+\sqrt{V_t}\,\dd W_t^{(1)},
\qquad
\dd V_t=\kappa(\theta-V_t)\dd t+\sigma\sqrt{V_t}\,\dd W_t^{(2)},
\qquad
\dd W^{(1)}\dd W^{(2)}=\rho\,\dd t.
\]
Posons $X_t=\ln S_t$. Par Itô (chapitre Itô) :
\[
\dd X_t=\frac{1}{S_t}\dd S_t-\frac12\frac{1}{S_t^2}(\dd S_t)^2
\]
or $(\dd S_t)^2=(\sqrt{V_t}S_t)^2(\dd W_t^{(1)})^2=V_t S_t^2\,\dd t$. Donc
\[
\boxed{\dd X_t=\Big(r-q-\tfrac12 V_t\Big)\dd t+\sqrt{V_t}\,\dd W_t^{(1)}.}
\]
Le couple $(X_t,V_t)$ est un processus de diffusion \emph{affine} (les coefficients sont au plus linéaires en $V$).

\subsection*{(2) La fonction caractéristique : définir l'objet et l'équation qu'il satisfait}
On cherche (pour $u\in\R$) la fonction caractéristique conditionnelle :
\[
\phi(t,x,v;u)=\E^{\Qmeas}\!\big[e^{iuX_T}\mid X_t=x,\ V_t=v\big].
\]
Par la théorie de Kolmogorov (ou Feynman--Kac), $\phi$ satisfait une PDE rétrograde :
\[
\partial_t \phi + \mathcal{L}\phi=0,
\qquad \phi(T,x,v;u)=e^{iux},
\]
où $\mathcal{L}$ est le générateur de $(X,V)$.

\subsection*{(3) Calcul du générateur $\mathcal{L}$ et ansatz affine}
À partir de
\[
\dd X=(r-q-\tfrac12 v)\dd t+\sqrt{v}\,\dd W^{(1)},
\qquad
\dd V=\kappa(\theta-v)\dd t+\sigma\sqrt{v}\,\dd W^{(2)},
\]
et $\dd W^{(1)}\dd W^{(2)}=\rho\,\dd t$, on obtient
\[
\mathcal{L}f
=(r-q-\tfrac12 v)\,\partial_x f+\kappa(\theta-v)\,\partial_v f
\tfrac12 v\,\partial_{xx}f+\tfrac12\sigma^2 v\,\partial_{vv}f+\rho\sigma v\,\partial_{xv}f.
\]
L'idée clé de Heston : chercher une solution de la forme affine
\[
\phi(t,x,v;u)=\exp\big(A(\tau;u)+B(\tau;u)\,v+iux\big),\qquad \tau=T-t.
\]
\paragraph{Dérivées.}
Si $\phi=\exp(A+Bv+iux)$, alors
\[
\partial_x\phi=iu\phi,\quad \partial_{xx}\phi=-u^2\phi,\quad
\partial_v\phi=B\phi,\quad \partial_{vv}\phi=B^2\phi,\quad
\partial_{xv}\phi=iuB\phi.
\]
\paragraph{Injection dans la PDE.}
On a $\partial_t\phi=-A'(\tau)\phi-B'(\tau)v\phi$ (car $\tau=T-t$).
En divisant la PDE par $\phi$ et en identifiant les coefficients constants et ceux en $v$, on obtient deux ODE :
\[
\boxed{
\begin{aligned}
B'(\tau)
&=\tfrac12\sigma^2 B(\tau)^2+(\rho\sigma iu-\kappa)B(\tau)-\tfrac12(u^2+iu),\\
A'(\tau)
&=iu(r-q)+\kappa\theta\,B(\tau),
\end{aligned}}
\qquad B(0)=0,\ A(0)=0.
\]
La première équation est une \textbf{Riccati} (non linéaire) ; la seconde est ensuite une intégration.

\subsection*{(4) Solution fermée de la Riccati : forme standard (à connaître)}
On pose
\[
d=\sqrt{(\kappa-\rho\sigma iu)^2+\sigma^2(u^2+iu)},
\qquad
g=\frac{\kappa-\rho\sigma iu-d}{\kappa-\rho\sigma iu+d}.
\]
Alors la solution s'écrit (forme «stable» utilisée en pratique) :
\[
\boxed{
B(\tau)
=\frac{\kappa-\rho\sigma iu-d}{\sigma^2}\cdot \frac{1-e^{-d\tau}}{1-ge^{-d\tau}}.
}
\]
Puis
\[
\boxed{
A(\tau)
=iu(r-q)\tau+\frac{\kappa\theta}{\sigma^2}\Big((\kappa-\rho\sigma iu-d)\tau-2\ln\frac{1-ge^{-d\tau}}{1-g}\Big).
}
\]
\paragraph{Remarque (branche complexe).}
Le choix de la racine carrée complexe et du logarithme est délicat numériquement (``little Heston trap'').
Le message M1 : la formule est semi-fermée, mais la stabilité dépend des conventions de branche.

\subsection*{(5) De la CF au prix : $P_1/P_2$ (structure) et intuition}
À $t=0$, $x=\ln S_0$ et $v=v_0$, donc
\[
\phi(u)=\exp\big(A(T;u)+B(T;u)v_0+iu\ln S_0\big).
\]
Le prix d'un call admet la structure :
\[
\boxed{
C=S_0e^{-qT}P_1-Ke^{-rT}P_2,
}
\]
où $P_1$ et $P_2$ sont calculés par des intégrales 1D impliquant $\phi$ (inversion de Fourier).
\paragraph{Intuition.}
La décomposition ressemble à BS :
\begin{itemize}[leftmargin=*]
\item $P_2$ joue le rôle d'une probabilité risque-neutre d'être ITM ;
\item $P_1$ est une probabilité sous une mesure «tiltée» (pondérée par $S_T$), d'où le terme $S_0e^{-qT}$.
\end{itemize}
Ce n'est pas qu'une analogie : on le démontre rigoureusement en manipulations de Fourier.

\subsection*{(6) Mini-exercice corrigé : retrouver BS comme cas limite}
\paragraph{Exercice.}
Montrer que si $V_t\equiv \bar v$ est constant (pas de vol-of-vol, pas de stochasticité de variance), alors $S$ suit un GBM de volatilité $\sqrt{\bar v}$ et on retrouve Black--Scholes.

\emph{Solution.}
Si $V_t\equiv\bar v$, la SDE devient
\[
\frac{\dd S_t}{S_t}=(r-q)\dd t+\sqrt{\bar v}\,\dd W_t,
\]
qui est exactement un GBM de volatilité $\sigma=\sqrt{\bar v}$.
Le pricing risque-neutre coïncide alors avec BS : la surface est plate (pas de skew), ce qui illustre que le skew vient de la stochasticité de $V$ et de la corrélation $\rho$.
